% 第 30 章　热力学第二、第三定律（完整正文）
\chapter{热力学第二、第三定律}

\step{反常识开场}

先来看一笔大到离谱的能源账。

一千克海水，温度每降低一度，会放出大约四千二百焦耳的热量。全球的海洋大约有一百亿亿亿（\(1.4\times10^{21}\)）千克水。让全部海水只降温一度，放出的热量约为 \(6\times10^{24}\) 焦耳——而全人类一年的总能耗大约是 \(6\times10^{20}\) 焦耳。也就是说，海水集体降温一度放出的热，够全人类按现在的规模用上约一万年。

于是有人设想了这样一艘船：它不需要煤、不需要油、不需要核燃料。它航行时从海水里抽取热量，用这些热发电驱动螺旋桨，再把温度稍微低了一点的海水排回大海。海水被抽走一点热，总量上毫发无损，下一艘船接着用。这艘船\textbf{不违反能量守恒}——它放出的每一份功，都有海水里减少的内能来付账，总账分毫不差。

这样的船，人类造出来了吗？一艘也没有。不是造不出来那么简单——是连一根能这样运转的螺丝钉都不存在。而且请注意，失败的方式很诡异：不是“效率太低”“成本太高”这种工程困难，而是它在原理上就连第一圈都转不起来。十九世纪的工程师试过各种方案，全部死在同一块暗礁上。

这种机器有个专门的名字，叫\textbf{第二类永动机}。它比第一类永动机（凭空造能的机器）狡猾得多：第一类永动机是赖账，账对不上，一眼就能抓；第二类永动机老老实实记账，分毫不差，却照样造不出来。抓它，需要一本完全不同的法律。

这块暗礁不在工程里，在定律里。能量守恒（第 28 章的主角）管得住“总量”，却管不住“方向”。本章要出场的，是热力学里专门管方向的两条定律：第二定律告诉你热为什么只能按一个方向用、热机的效率为什么有铁的上限；第三定律告诉你温度的地板上为什么永远站不上去。

\step{先猜一猜}

在往下读之前，先写下你对下面这个问题的猜测：那艘“烧海水”的船为什么造不出来？

第一种猜测：是技术不够好。就像古人造不出飞机，将来的工程师也许能造出来，只是我们现在还不会。

第二种猜测：海水的热虽然多，但太“稀”了——温度太低的热不好用，只要找到把低温水里的热“浓缩”成高温热的办法，船就能跑。

第三种猜测：存在一条和能量守恒彼此独立的新定律，它专门禁止“只从单一热源取热做功”，无论技术怎么进步都绕不过去。

顺带再猜一个家里的问题：夏天厨房太热，有人提议把冰箱门打开当空调用。你觉得房间会变凉、变热，还是不变？同样先写下你的判断。这一章结束时，我们要回来对两次答案。

\step{动手实验}

\begin{experiment}[开门的冰箱能当空调吗]
你需要：一台正常工作的冰箱、一支温度计（厨房的、体温计不行，要能测室温的）、一部手机用来计时。

第一步：把温度计放在厨房里离冰箱两米以外的地方，等五分钟，记下室温。

第二步：打开冰箱门，让它敞着。为了让冰箱持续制冷，可以用一卷胶带把门框上感应关门的小开关压住（有些冰箱没有，那就不用管），让压缩机一直工作。等一到两个小时。期间别开大功率电器。

第三步：再读一次温度计。同时做两个附加观察：摸一摸冰箱背面或两侧，那里是热的还是凉的？再看看冷藏室内部，是不是真的在变凉？

你会看到：冷藏室里确实更冷了，但房间的平均温度不降，反而微微上升。冰箱辛苦工作一两个小时，把厨房变得更热了。
\end{experiment}

这个结果值得停三秒钟。冰箱明明在“制冷”，为什么房间反而升温？拆开看冰箱在做什么：它从冷藏室（低温处）把热抽出来，通过背后的散热管排到房间（高温处）里——热在从冷的地方流向热的地方！这听起来已经像是“反向”了。但注意代价：压缩机一直在耗电，而电能最终也全部变成热，排进同一个房间。房间收到的热，等于从冷藏室搬来的热\textbf{加上}电费买来的热。两笔都是正的，房间当然变热。

真正的空调为什么能给房间降温？因为它的散热部分（外机）挂在窗外，把热排给了室外。同一个道理：想给屋子降温，必须有一个“屋子外面”可以排热。这个朴素的观察——\textbf{把热从冷处搬到热处，必须额外付功；而且总有一个更大的环境在接收账单}——就是本章第一条主角定律的日常模样。

这个实验朴素，但请先写下三个观察题的答案，本章的概念会逐个接住它们：冰箱后背的热，是从冷藏室“搬”出来的，还是压缩机“造”出来的？如果房间是严格隔热的，开着冰箱门一直等下去，房间升温会停下来吗？把温度计伸进冷藏室，那里的确在变冷——一个让局部变冷、让整体变热的机器，和我们这章要找的“方向”有什么关系？

\step{旧模型遇到麻烦}

回到第 28 章。热力学第一定律说：内能的变化等于吸收的热与外界做功的总账，能量不生不灭。这条定律威力巨大，它一票否决了所有“凭空造能”的机器——那种机器叫第一类永动机，死得不冤。

但请用第一定律审判下面三个过程，你会发现它一次也判不下去：

凉咖啡从房间里慢慢吸热、自己重新沸腾——账是平的：房间少的热正好等于咖啡多的热。冰水里的一小块冰变得更冰、周围的水变得更烫——账也是平的：水失去的热正好等于冰得到的热（方向反了而已）。烧海水的船——账还是平的，开场已经算过。

三个过程全都满足能量守恒，三个过程全都从未被任何人见过。第一定律是一位只核对总金额、从不过问资金流向的会计。我们需要一条新定律来回答：\textbf{在能量守恒允许的所有过程里，哪些真的会发生？}

历史上有意思的是，这条新定律的第一块基石，是在能量守恒还没建立的时候奠定的。1824 年，二十八岁的法国工程师萨迪·卡诺出版了一本小册子，问了一个工程师式的问题：蒸汽机的效率有没有上限？当时的蒸汽机效率只有可怜的百分之几，工程师们不知道剩下的百分之九十多是“还没改进到位”，还是“本来就不许拿”。卡诺用一个理想化的思想实验给出了答案：\textbf{有上限，而且这个上限只取决于热机工作时的两个温度，跟用什么工质、什么结构毫无关系}。讽刺的是，卡诺当时还相信错误的“热质说”，却推出了正确的结论——好的思想实验有时比错的理论框架更有生命力。二十多年后，克劳修斯和开尔文把卡诺的洞见重新安放在能量守恒的地基上，这就是热力学第二定律。

\begin{mainline}
旧模型的麻烦可以压成一句话：\textbf{能量守恒管得住账本的总额，管不住每一笔账的流向。}三个“合法却从未发生”的过程，方向全部一致——热总是自发地从高温流向低温，从不自动回头。新定律必须从“方向”入手。
\end{mainline}

\step{发明新概念}

历史给我们留下了第二定律的两种经典表述，它们像同一枚硬币的两面。

\textbf{克劳修斯表述}：不可能把热量从低温物体传到高温物体，而不引起其他变化。注意那个尾巴——“而不引起其他变化”。冰箱明明在把热从冷处搬到热处，它违反了吗？没有，因为它引起了“其他变化”：电表在转，电能在变成热。克劳修斯表述禁止的是\textbf{不付代价}的逆流。

\textbf{开尔文表述}：不可能从单一热源吸收热量，使之完全变为功，而不产生其他影响。“单一热源”四个字正中那艘烧海水的船：海水就是它唯一的热源，船上没有任何更冷的地方可以排热，于是吸来的热一分也变不成功的净输出。注意它同样没有说“热不能变功”——热机天天在变功；它说的是\textbf{不能全部变}，也不能\textbf{只从一个热源变}。

这两种表述其实是同一条定律。可以做一个简短的推理：假设有人造出了违反开尔文表述的机器，能从单一热源取热全部变功；把这台机器的功拿去驱动一台普通冰箱，从冷库里抽热排到那个热源——净效果就是热不付任何代价地从冷处流到了热处，克劳修斯表述跟着垮掉。反过来，违反克劳修斯表述的机器配上普通热机，也能拼出开尔文禁止的东西。两条表述要么同生，要么同死。

那么卡诺说的“效率上限”是怎么回事？卡诺设计了一台存在于思想里的完美热机：它只和两个恒定温度的热库打交道——一个高温热库（温度 \(T_1\)），一个低温热库（温度 \(T_2\)）；它工作时先接触高温库吸热，再与热库脱开、膨胀做功降温，然后接触低温库排热，最后被压缩回出发点。整个过程慢到每一步都近乎平衡，没有摩擦、没有温差乱流，随时可以让它掉头倒着走。这样的循环叫\textbf{卡诺循环}，这样的热机叫可逆热机。卡诺证明：在同样两个热库之间，一切可逆热机的效率都相同；任何实际（不可逆）热机的效率都不可能超过它。而可逆机的效率只由两个温度决定——这甚至给了人类一把不依赖任何具体物质的温度计：测两台热机的效率比，就能定义温度的比值。这把“思想温度计”定义出来的温标，就是绝对温标，单位开尔文。

\begin{center}
\begin{tikzpicture}[scale=1.0, >={Stealth[length=2.5mm]}]
  % 高温热库
  \draw[thick, fill=red!12] (-2.2,3.4) rectangle (2.2,4.2);
  \node at (0,3.8) {高温热库 \(T_1\)};
  % 低温热库
  \draw[thick, fill=blue!10] (-2.2,-1.0) rectangle (2.2,-0.2);
  \node at (0,-0.6) {低温热库 \(T_2\)};
  % 热机
  \draw[thick, fill=yellow!15] (0,1.6) circle (0.85);
  \node at (0,1.6) {热机};
  % Q1 箭头
  \draw[->, very thick, red!60!black] (-0.55,3.4) -- (-0.55,2.5);
  \node[left] at (-0.6,2.95) {吸热 \(Q_1\)};
  % Q2 箭头
  \draw[->, very thick, blue!60!black] (0.55,0.7) -- (0.55,-0.2);
  \node[right] at (0.62,0.25) {排热 \(Q_2\)};
  % W 箭头
  \draw[->, very thick] (0.85,1.6) -- (2.1,1.6);
  \node[right] at (2.15,1.6) {输出功 \(W=Q_1-Q_2\)};
\end{tikzpicture}
\end{center}

这张图是本章最重要的画面，值得盯着看十秒：热机\textbf{不是把热“变”成功的转换器，而是让热从高温流向低温、在半路截下一部分变功的水轮机}。截下多少是功（\(W\)），放走多少是废热（\(Q_2\)），比例由两个温度说了算。废热不是机器不够好造成的浪费，而是机器运转的必要条件——没有下游，就没有水流。

用第 29 章的语言，第二定律还有一个更根本的表述：孤立系统的熵几乎总是增加，平衡时达到最大，记作 \(\Delta S \geq 0\)。热从高温流向低温为什么熵增？同样一份热量 \(Q\)，放在高温处时“分摊”到每度温度上较少，放在低温处时分摊较多——所以热从高温挪到低温，总账是赚的。下一节我们把这个直觉变成公式。

\step{一句话模型}

\begin{oneline}
热机是让热从高温流向低温、在半路截一部分变功的机器；截留的比例有铁的上限——只由两个温度决定；所以单一热源的热一分也变不成功，而越接近绝对零度，剩下的热越排不出去，零度本身只能逼近、永远到不了。
\end{oneline}

这句话里藏着三个“永远”：废热永远排掉一部分（开尔文），逆流永远要付代价（克劳修斯），零度永远差一步（第三定律）。它们不是三条禁令，是同一条统计规律（第 29 章的熵增）在热机、制冷机和温度计上的三张面孔。

\step{数学翻译器}

现在给思想装上数字。设可逆热机从高温热库 \(T_1\) 吸热 \(Q_1\)，向低温热库 \(T_2\) 排热 \(Q_2\)，对外做功 \(W=Q_1-Q_2\)。效率定义为你拿到的功占你投入的热的比例：
\[
\eta=\frac{W}{Q_1}=\frac{Q_1-Q_2}{Q_1}=1-\frac{Q_2}{Q_1}.
\]

可逆意味着整个循环不产生熵：高温库交出热量 \(Q_1\)，失去熵 \(Q_1/T_1\)；低温库收到热量 \(Q_2\)，得到熵 \(Q_2/T_2\)。可逆要求两者恰好相抵：
\[
\frac{Q_1}{T_1}=\frac{Q_2}{T_2}
\quad\Longrightarrow\quad
\frac{Q_2}{Q_1}=\frac{T_2}{T_1}.
\]
代回效率公式，得到卡诺上限：
\[
\boxed{\ \eta \leq 1-\frac{T_2}{T_1}\ }
\]
注意这里的温度必须是绝对温标（开尔文），从绝对零度算起——摄氏度的“零”是人为规定的冰点，进不了这个公式。

公式写出来，立刻用特殊情况检查它。第一，\(T_2=T_1\)：没有温差，\(\eta=0\)——两个同温度的热库之间开不出热机，正如上下游同高的河上建不了水电站，烧海水的船再次沉没。第二，\(T_2\to 0\)：\(\eta\to 100\%\)，但等会儿第三定律会告诉你 \(T_2=0\) 这个热库不存在，所以效率永远差一口气到不了百分之百——两个定律在这里扣上了环。第三，方向反转：如果 \(T_2>T_1\)，公式给出负效率，荒唐吗？不荒唐——那正是你把机器倒过来开的情形：外界输入功，把热从冷库搬到热库，热机变成了制冷机。公式的“翻脸”恰好对应机器的“倒车”，说明式子写对了。

下面这张图把卡诺循环画在压强—体积平面上（第 28 章说过，这个图上一条闭合曲线围成的面积就是循环对外做的净功）：

\begin{center}
\begin{tikzpicture}[scale=1.15, >={Stealth[length=2.2mm]}]
  \draw[->] (0,0) -- (5.6,0) node[right] {体积 \(V\)};
  \draw[->] (0,0) -- (0,3.9) node[above] {压强 \(p\)};
  % 等温吸热 T1: A->B
  \draw[thick, red!60!black] plot[smooth, tension=0.9] coordinates {(1.0,3.3) (1.6,2.7) (2.4,2.15) (3.2,1.75)};
  % 绝热膨胀: B->C
  \draw[thick] plot[smooth, tension=0.9] coordinates {(3.2,1.75) (3.9,1.15) (4.6,0.72)};
  % 等温放热 T2: C->D
  \draw[thick, blue!60!black] plot[smooth, tension=0.9] coordinates {(4.6,0.72) (3.2,0.85) (1.7,1.15)};
  % 绝热压缩: D->A
  \draw[thick] plot[smooth, tension=0.9] coordinates {(1.7,1.15) (1.3,2.1) (1.0,3.3)};
  \fill (1.0,3.3) circle (0.05) node[above left] {A};
  \fill (3.2,1.75) circle (0.05) node[above right] {B};
  \fill (4.6,0.72) circle (0.05) node[below right] {C};
  \fill (1.7,1.15) circle (0.05) node[below left] {D};
  \node[red!60!black] at (2.0,3.35) {等温吸热 \(Q_1\)（\(T_1\)）};
  \node[blue!60!black] at (3.6,0.35) {等温放热 \(Q_2\)（\(T_2\)）};
  \node[rotate=-32] at (4.35,1.45) {\small 绝热膨胀};
  \node[rotate=62] at (0.88,2.1) {\small 绝热压缩};
\end{tikzpicture}
\end{center}

A 到 B：贴着高温库慢慢膨胀，吸热 \(Q_1\)；B 到 C：脱离热库继续膨胀，温度自己降到 \(T_2\)；C 到 D：贴着低温库被压缩，放热 \(Q_2\)；D 到 A：脱离热库继续压缩，温度回升到 \(T_1\)，循环闭合。注意 B、C 和 D、A 这两段不与任何热库交换热量（绝热段）——整个机器只在两个温度上与外界“说话”，这正是它可逆的秘密：任何时候让它倒转，每一步都能原路退回。

\begin{mathbridge}
上面用“可逆则熵不变”一步跳到了卡诺效率，其实也能从理想气体一步步推出来，顺带看一眼温度是怎样混进效率的。等温过程中理想气体内能不变，吸的热全变成膨胀功：\(Q_1=nRT_1\ln(V_B/V_A)\)，同理 \(Q_2=nRT_2\ln(V_C/V_D)\)。两段绝热过程满足 \(T V^{\gamma-1}=\text{常量}\)，把它们写出来一比，恰好给出 \(V_B/V_A=V_C/V_D\)。于是两个对数相等，相除即得 \(Q_2/Q_1=T_2/T_1\)。推导里气体的一切性质（\(n\)、\(\gamma\)、体积）全部约掉，只剩温度比——这就是“效率与工质无关”的代数证据。

制冷机也有对应的上限。把卡诺循环倒着跑，输入功 \(W\)，从冷库 \(T_2\) 抽热 \(Q_2\)，制冷系数
\[
\mathrm{COP}_{\text{冷}}=\frac{Q_2}{W}\leq\frac{T_2}{T_1-T_2}.
\]
拿这个公式检查绝对零度：固定房间温度 \(T_1\)，让冷库温度 \(T_2\) 越来越低，COP 一路跌向零——从越冷的地方搬同样一份热，要付的功越多；逼近零度时，搬走一丁点热都要付出趋于无穷的功。\textbf{第三定律的“零度不可达”，早就写在这行公式里。}
\end{mathbridge}

现在做两个带真实数字的估算。

第一笔账，火力发电厂。现代超超临界机组的蒸汽温度约为 \(600\,^\circ\mathrm{C}\)，合 \(873\ \mathrm{K}\)；冷凝器用江水或冷却塔散热，约 \(30\,^\circ\mathrm{C}\)，合 \(303\ \mathrm{K}\)。卡诺上限：
\[
\eta\leq 1-\frac{303}{873}\approx 65\%.
\]
实际最好的机组约为 \(45\%\)——材料耐温、传热温差、摩擦各吃掉一块。卡诺没有骗我们：他说的是“天花板”，不是“保证书”。但你倒过来用它就有力量了：谁要是宣称造出了烧 \(600\,^\circ\mathrm{C}\) 蒸汽、排 \(30\,^\circ\mathrm{C}\) 环境、效率七成的机器，不用看图纸，直接判定违反第二定律。

第二笔账，地热。某地热田喷出 \(150\,^\circ\mathrm{C}\)（\(423\ \mathrm{K}\)）的热水，环境 \(25\,^\circ\mathrm{C}\)（\(298\ \mathrm{K}\)）：上限 \(1-298/423\approx 30\%\)，实际电站往往只有百分之十几。温度不高的能源，先天就戴着低的效率天花板——这就是为什么“低温余热发电”听起来美好、算完账常常不划算。

最后把温度的地板请出来。\textbf{热力学第三定律}有两个等价的说法。其一（能斯脱表述）：当温度趋近绝对零度时，完美晶体的熵趋近于零——用第 29 章的话说，零度时系统只剩一个微观状态（\(\Omega\to1\)），房间只剩一间，没得更“大”可去。其二（不可达表述）：不可能用有限步骤的操作把任何系统冷却到绝对零度。直觉是这样：降温的本质是把热排给更冷的东西，而“最冷”没有更冷者来接盘；越接近零度，可排的热越少、传热越慢、制冷系数越小（上面 mathbridge 的公式），每降一步都比上一步更难，台阶无穷多，而你的步数有限。实验史就是一部逼近史：液氦能把东西冷到 \(4.2\ \mathrm{K}\)，稀释制冷机到约 \(0.01\ \mathrm{K}\)，激光冷却原子云到十亿分之一开量级——纪录不断刷新，零度始终还差最后一步，而且永远会差。

\step{模型的边界}

老规矩，交代三件事：假设了什么、在哪里成立、在哪里会被误用。

\textbf{第二定律是统计定律，不是铁律。}第 29 章已经见过：熵增是“几乎必然”，不是“绝对禁止”。在足够小的系统里，涨落可以短时间、小尺度地“逆流”——悬浮在水中的微米颗粒会被分子撞得短暂获得一点热运动的功。这\textbf{不违反}第二定律，因为它说的本来就是宏观平均。但请别激动：想靠收集涨落造一台稳定的纳米热机是不行的，因为只要你要求它\textbf{持续、可靠}地输出功，平均效果就必然回到熵增一边。“几乎不可能”与“绝对禁止”的区别，正是第二定律最深刻的性格：它管辖的是大数世界。

\textbf{卡诺上限是判据，不是工程承诺。}它假设：只有两个恒温热库、过程无限缓慢（准静态）、无摩擦无泄漏、完全可逆。实际热机为了提高功率，必须让传热有足够温差，而温差传热本身就产生熵——真实世界里“快”和“省”永远在打架。正确用法是拿上限当照妖镜（超过上限的宣称必假）和当改进的路标（离上限还远，说明有空间）。

\textbf{别把“热”和“温度”混进日常口号。}“把环境里的热全部变成功”之所以不可能，要害不在热少，而在没有更冷的地方排熵。反过来说，冬天给房间供暖，最聪明的办法不是把电“烧成”热（电暖器，一份电一份热），而是用一份电去\textbf{搬}好几份热（热泵）。后者听起来像“一份投入三份产出”，会让直觉报警——它不违反能量守恒，因为多出来的热是从室外搬来的，电只付了搬运费。这个反直觉的例子是第二定律送我们的礼物，不是漏洞。

\begin{challenge}
还有一个更隐蔽的例外区：自引力系统。卫星在大气阻力下损失能量、轨道降低，速度反而变快；星团收缩时“温度”反而升高——引力系统的“热容”是负的，第二定律的通常表述在引力主导的体系里要重新小心陈述。更惊人的是黑洞：黑洞不但有熵（正比于表面积），还有温度，会极其缓慢地向外辐射。热力学三定律在黑洞那里几乎逐字成立，只是账本上多了引力这一栏。这条线一直通向广义相对论和量子引力的未解之地——本书后面还会回来敲门。
\end{challenge}

\begin{mainline}
使用说明：宏观、近平衡、有明确热库时，第二定律的方向判断和卡诺上限可以放心使用；系统很小要看涨落，引力很强要重写规则；把卡诺上限当实际效率、把统计规律当绝对禁令、把热泵“超百分之百”当永动机，是本章三大经典误用。
\end{mainline}

\step{回头解释开场现象}

现在审判那艘烧海水的船。

它唯一的热库是海水。要从海水吸热并输出净功，按卡诺的框架，必须同时存在一个更冷的地方接收废热——船上有吗？没有。周围的海水、空气和船身在同一个温度附近。没有下游，水轮机空转都不转：这正是开尔文表述禁止的“从单一热源取热全部变功”。换成熵的语言更干脆：海水失热降温，熵减少；排出的废热为零，没有谁的熵增加来补账；孤立系统总熵要下降——概率上被判了死刑，而且不是“很难”，是“零的个数排到太阳”那种不可能（第 29 章）。

再补一刀：假如这艘船真的存在，把它的输出功接上一台普通冰箱，就能让热从冷海水自动流进热空气——克劳修斯表述跟着崩塌，凉咖啡可以自己沸腾，空调不用电。一条船沉下去，会拖垮整条定律的舰队；反过来，定律站住了，船就永远造不出。这不是技术问题，是方向问题。

开门冰箱的答案现在也是明牌：冰箱从冷藏室往房间搬热，这一步本来就在“逆流”，所以它必须耗电付代价；而它付的电费最终也变成热留在房间里。房间收到的热等于“搬来的热加电费的热”，两笔都为正，所以房间只会变热。你实验里摸到的冰箱后背，就是这台机器在向你展示克劳修斯表述的收据。真正的空调能降温，是因为它的“房间外面”存在——外机把热排给了整个大气。

开头的三个猜测，现在对答案。猜测一（技术不够）错：再好的技术也不能让定律让路，正如再精密的杠杆也撬不动动量守恒。猜测二（把低温热浓缩）错：浓缩热本身就是把热从冷处搬到热处，每一步都要付功，付的功比赚的多——卡诺上限保证了这笔买卖永远亏损。猜测三对：热力学第二定律独立于能量守恒，它不管总量、专管方向。

\step{脑洞实验与挑战题}

\begin{thought}[把太阳当锅炉，把银河当冷凝器]
来设计宇宙尺度的热机。锅炉用太阳表面，约 \(5800\ \mathrm{K}\)；冷凝器用整个宇宙——深空里弥漫着宇宙微波背景辐射，温度只有 \(2.7\ \mathrm{K}\)，它是任何辐射最终可以排热的“大冷坑”。卡诺上限：
\[
\eta\leq 1-\frac{2.7}{5800}\approx 99.95\%.
\]
几乎完美。这就是“用星光发电”的终极效率天花板。但注意最后一块骨头：你的冷凝器再冷也有 \(2.7\ \mathrm{K}\)，所以废热必须以高于 \(2.7\ \mathrm{K}\) 的温度排给深空，那 \(0.05\%\) 永远拿不回来；而宇宙背景本身还在随膨胀缓慢降温，宇宙的“冷坑”在变冷，却永远不会是零。热机的效率上限、制冷机的零度下限，在同一片星空下握手。
\end{thought}

\begin{thought}[冰箱和生命违反第二定律吗]
冰在冰箱里结成整齐的晶体，受精卵长成结构精密的人——局部的熵都在大幅下降。第二定律出错了吗？你已经有了全部工具：检查系统边界。冰箱内部不是孤立系统，它排给厨房的热更多；人不是孤立系统，一辈子吃的、呼吸的、散掉的热加起来，让“人加环境”的总熵只增不减。局部有序是用更大范围的无序买来的。这个主题太大，第 32 章会专门展开，届时麦克斯韦妖也会登场收账。
\end{thought}

\begin{puzzles}
\item \textbf{热泵的“超百分之百”。}广告说：某热泵冬天“消耗一度电，给室内送来四度电的热”，制热效率 \(400\%\)。电暖器最多 \(100\%\)。请问：这违反能量守恒吗？为什么严寒天（室外零下十几度）热泵的这项指标会明显下滑？

  \textit{思路提示：} 不违反。电暖器是把电“变成”热，一份变一份；热泵是用电“搬”热，从室外冷空气中抽取热量排进室内，多出来的三份来自室外，不是凭空产生。卡诺热泵的供热系数上限为 \(T_1/(T_1-T_2)\)：室内外温差越大，分母越大，搬运越费力，指标自然下滑——公式里早写着天气的脾气。

\item \textbf{一杯水的温差阴谋。}有人设想：把一杯 \(50\,^\circ\mathrm{C}\) 的水插上一块隔板，设法让左半杯降到 \(40\,^\circ\mathrm{C}\)、右半杯升到 \(60\,^\circ\mathrm{C}\)，然后用两半的温差驱动一台小热机发电，发完电水又回到 \(50\,^\circ\mathrm{C}\)，循环往复。能量看起来完全守恒。请指出这个方案的致命一步在哪里。

  \textit{思路提示：} 致命一步不在热机，在“让一半变冷、另一半变热”——这是热自发从热处流向冷处的逆过程，等同于不付代价地制造温差，直接违反克劳修斯表述。后面那台小热机是清白的，但它等不到它的上游。判断永动机时，盯住每一步的方向，而不只是总账。

\item \textbf{零度接力赛。}有人说：虽然一台制冷机到不了绝对零度，但可以用接力——先用甲机冷到 \(1\ \mathrm{K}\)，再用乙机接力到 \(10^{-3}\ \mathrm{K}\)，再用丙机……每一级都更近一步，为什么不行的？

  \textit{思路提示：} 每一级制冷都需要一个更冷的热库来接废热，而“更冷的热库”正是你要造的东西本身；同时制冷系数正比于 \(T_2/(T_1-T_2)\)，越接近零度，搬走同样的热需要的功越趋近无穷。台阶有无穷多级，而任何现实操作只有有限步——这正是第三定律“不可达”表述的内容。逼近可以，到达不行。

\item \textbf{尾气回收的边界。}汽车尾气带走大量热量，有人提议：“把尾气的热全部回收，变成动力，效率就能到百分之百。”这句话错在哪一层？如果只能回收一部分，上限由什么决定？

  \textit{思路提示：} 回收尾气热也是热机，必须有冷端——最终还得向环境（约 \(300\ \mathrm{K}\)）排热。尾气温度就算有 \(700\ \mathrm{K}\)，回收环节的卡诺上限也只有 \(1-300/700\approx 57\%\)。“全部回收”等价于从尾气这个高温库取热全部变功、不向任何更冷处排热，直接撞上开尔文表述。
\end{puzzles}

本章从一艘永远造不出的船出发，找到了热世界的交通规则：能量守恒管总额，第二定律管方向，卡诺用两个温度给热机的胃口定了铁的上限，第三定律则封死了温度的地板。至此，热学三定律齐了——第零定律给了温度存在的资格，第一定律给了账本，第二定律给了方向，第三定律给了边界。下一章我们要换一副眼镜：不再只看宏观的热库和热机，而是钻进分子的亿万大军，看压强、温度、熵这些宏观概念是怎样从统计里“长”出来的。那条路，会把第 29 章数微观状态的手艺，发展成一整套统计物理。
